% =========================================================
% Approximation Theorems — Notes
% Chapter: continuity
% Section: approximation
% Volume III · Analysis
% =========================================================

\section{Approximation of Continuous Functions}
\label{sec:approximation}

% ---------------------------------------------------------
% Toolkit
% ---------------------------------------------------------

\begin{tcolorbox}[title={Toolkit — Approximation}]
\begin{itemize}
  \item Approximation hierarchy: step functions $\subset$ piecewise
    linear functions $\subset$ polynomials. Each class is dense in
    $C([a,b])$ under the uniform norm.
  \item All three results use Heine--Cantor: uniform continuity on
    $[a,b]$ produces the $\delta$ that controls the partition.
  \item Bernstein polynomials give a constructive proof of Weierstrass:
    $B_n(x) = \sum_{k=0}^n f(k/n)\binom{n}{k}x^k(1-x)^{n-k}$.
  \item Step function approximation feeds directly into Riemann
    integrability of continuous functions.
\end{itemize}
\end{tcolorbox}

% ---------------------------------------------------------
% Theorem — Step Function Approximation
% ---------------------------------------------------------

\begin{theorem}[Step Function Approximation]
\label{thm:step-function-approximation}
Let $I = [a,b]$ be a closed bounded interval and let
$f : I \to \mathbb{R}$ be continuous on $I$. For every $\varepsilon
> 0$ there exists a step function $s_\varepsilon : I \to \mathbb{R}$
such that
\[
  |f(x) - s_\varepsilon(x)| < \varepsilon
  \qquad\text{for all } x \in I.
\]

\smallskip
\noindent
\hyperref[prf:step-function-approximation]{\textit{Go to proof.}}
\end{theorem}

\begin{remark*}[Standard quantified statement]
\[
  f \text{ cts on } [a,b],\;
  \varepsilon > 0
  \;\Rightarrow\;
  \exists \text{ step function } s_\varepsilon:\;
  \sup_{x \in [a,b]} |f(x) - s_\varepsilon(x)| < \varepsilon.
\]
\end{remark*}

\begin{remark*}[Interpretation]
Heine--Cantor provides $\delta > 0$ such that $|f(x) - f(y)|
< \varepsilon$ whenever $|x-y| < \delta$. Partition $[a,b]$ into
subintervals of width less than $\delta$ and define $s_\varepsilon$
to be constant on each piece at the value of $f$ at the left
endpoint. Uniform continuity ensures the approximation error is
less than $\varepsilon$ throughout. This is used in the theory of
Riemann integration to show that continuous functions on $[a,b]$
are Riemann integrable.
\end{remark*}

% ---------------------------------------------------------
% Corollary — Uniform Partition Version
% ---------------------------------------------------------

\begin{corollary}[Uniform Partition Approximation]
\label{cor:uniform-partition-approximation}
Let $I = [a,b]$ and let $f : I \to \mathbb{R}$ be continuous on $I$.
For every $\varepsilon > 0$ there exists $m \in \mathbb{N}$ such that
if $I$ is divided into $m$ equal subintervals of length $h =
(b-a)/m$, the step function $s_\varepsilon$ defined by
\[
  s_\varepsilon(x) = f(x_{k-1}),
  \quad x \in [x_{k-1}, x_k],
\]
satisfies $|f(x) - s_\varepsilon(x)| < \varepsilon$ for all $x \in I$.
\end{corollary}

\begin{remark*}[Standard quantified statement]
\[
  \forall \varepsilon > 0,\;
  \exists m \in \mathbb{N}:\;
  \sup_{x \in [a,b]} |f(x) - s_\varepsilon(x)| < \varepsilon,
\]
where $s_\varepsilon$ is the uniform left-endpoint step function
with $m$ equal subintervals.
\end{remark*}

\begin{remark*}[Interpretation]
This is the uniform partition specialization of the Step Function
Approximation Theorem. The Heine--Cantor $\delta$ determines $m$
via $m > (b-a)/\delta$. This is the form directly relevant to
Riemann sums: upper and lower sums over uniform partitions converge
to the integral for continuous functions.
\end{remark*}

% ---------------------------------------------------------
% Definition — Piecewise Linear Function
% ---------------------------------------------------------

\begin{definition}[Piecewise Linear Function]
\label{def:piecewise-linear-function}
Let $I = [a,b]$. A function $g : I \to \mathbb{R}$ is
\textbf{piecewise linear} on $I$ if $I$ is the union of finitely
many disjoint subintervals $I_1, \ldots, I_m$ such that the
restriction of $g$ to each $I_k$ is an affine (linear) function.
\end{definition}

\begin{remark*}[Standard quantified statement]
\[
  g \text{ piecewise linear on } [a,b]
  \;\iff\;
  \exists a = x_0 < x_1 < \cdots < x_m = b,\;
  \forall k:\;
  g\big|_{[x_{k-1},x_k]} \text{ is affine.}
\]
\end{remark*}

\begin{remark*}[Interpretation]
A piecewise linear function is continuous (the affine pieces are
required to agree at the partition points) and its graph consists
of finitely many line segments. It is more regular than a step
function — it has no jump discontinuities — but less regular than
a polynomial.
\end{remark*}

% ---------------------------------------------------------
% Theorem — Piecewise Linear Approximation
% ---------------------------------------------------------

\begin{theorem}[Piecewise Linear Approximation]
\label{thm:piecewise-linear-approximation}
Let $I = [a,b]$ and let $f : I \to \mathbb{R}$ be continuous on $I$.
For every $\varepsilon > 0$ there exists a continuous piecewise
linear function $\ell_\varepsilon : I \to \mathbb{R}$ such that
\[
  |f(x) - \ell_\varepsilon(x)| < \varepsilon
  \qquad\text{for all } x \in I.
\]

\smallskip
\noindent
\hyperref[prf:piecewise-linear-approximation]{\textit{Go to proof.}}
\end{theorem}

\begin{remark*}[Standard quantified statement]
\[
  f \text{ cts on } [a,b],\;
  \varepsilon > 0
  \;\Rightarrow\;
  \exists \text{ cts piecewise linear } \ell_\varepsilon:\;
  \sup_{x \in [a,b]} |f(x) - \ell_\varepsilon(x)| < \varepsilon.
\]
\end{remark*}

\begin{remark*}[Interpretation]
Heine--Cantor gives $\delta$; partition $[a,b]$ into subintervals of
width less than $\delta$ at nodes $x_0 < \cdots < x_n$; define
$\ell_\varepsilon$ on each $[x_{j-1}, x_j]$ as the affine function
through $(x_{j-1}, f(x_{j-1}))$ and $(x_j, f(x_j))$. For any
$x \in [x_{j-1}, x_j]$, the value $\ell_\varepsilon(x)$ is a convex
combination of $f(x_{j-1})$ and $f(x_j)$, which by the IVT equals
some value of $f$ on the same subinterval, so within $\varepsilon$
of $f(x)$.
\end{remark*}

% ---------------------------------------------------------
% Theorem — Weierstrass Approximation Theorem
% ---------------------------------------------------------

\begin{theorem}[Weierstrass Approximation Theorem]
\label{thm:weierstrass-approximation}
Let $I = [a,b]$ and let $f : I \to \mathbb{R}$ be continuous on $I$.
For every $\varepsilon > 0$ there exists a polynomial $p_\varepsilon$
such that
\[
  |f(x) - p_\varepsilon(x)| < \varepsilon
  \qquad\text{for all } x \in I.
\]

\smallskip
\noindent
\hyperref[prf:bernstein-approximation]{\textit{Go to proof.}}
\end{theorem}

\begin{remark*}[Standard quantified statement]
\[
  f \text{ cts on } [a,b],\;
  \varepsilon > 0
  \;\Rightarrow\;
  \exists p_\varepsilon \in \mathbb{R}[x]:\;
  \sup_{x \in [a,b]} |f(x) - p_\varepsilon(x)| < \varepsilon.
\]
\end{remark*}

\begin{remark*}[Interpretation]
Every continuous function on a closed bounded interval is the
uniform limit of polynomials. In the language of function spaces,
the polynomials are dense in $C([a,b])$ under the uniform norm
$\|f\|_\infty = \sup_{x \in [a,b]} |f(x)|$. This is a major result:
polynomials, which are the simplest functions, are sufficient to
approximate any continuous function to arbitrary precision. The
constructive proof via Bernstein polynomials is given below.
\end{remark*}

% ---------------------------------------------------------
% Definition — Bernstein Polynomial
% ---------------------------------------------------------

\begin{tcolorbox}[title={Definition --- Bernstein Polynomial}]
\begin{definition}[Bernstein Polynomial]
\label{def:bernstein-polynomial}
Let $f : [0,1] \to \mathbb{R}$. The \textbf{$n$-th Bernstein
polynomial for $f$} is
\[
  B_n(x) \;:=\; \sum_{k=0}^{n}
  f\!\left(\tfrac{k}{n}\right)
  \binom{n}{k} x^k (1-x)^{n-k},
  \qquad x \in [0,1].
\]
\end{definition}
\end{tcolorbox}

\begin{remark*}[Standard quantified statement]
\[
  B_n(x) = \sum_{k=0}^{n}
  f\!\Bigl(\frac{k}{n}\Bigr)
  \binom{n}{k} x^k (1-x)^{n-k},
  \quad
  \binom{n}{k} = \frac{n!}{k!(n-k)!}.
\]
\end{remark*}

\begin{remark*}[Interpretation]
$B_n(x)$ is a polynomial of degree at most $n$. Its coefficients
are weighted values of $f$ at the equally spaced points $0, 1/n,
2/n, \ldots, 1$, weighted by the binomial probabilities
$\binom{n}{k}x^k(1-x)^{n-k}$. The probabilistic interpretation
is direct: $B_n(x) = \mathbb{E}[f(S_n/n)]$ where $S_n \sim
\mathrm{Binomial}(n,x)$. The law of large numbers forces $S_n/n
\to x$ in probability, and uniform continuity of $f$ converts
this convergence into uniform approximation.
\end{remark*}

% ---------------------------------------------------------
% Theorem — Bernstein Approximation Theorem
% ---------------------------------------------------------

\begin{theorem}[Bernstein Approximation Theorem]
\label{thm:bernstein-approximation}
Let $f : [0,1] \to \mathbb{R}$ be continuous and let
$\varepsilon > 0$. There exists $n_\varepsilon \in \mathbb{N}$
such that for all $n \geq n_\varepsilon$,
\[
  |f(x) - B_n(x)| < \varepsilon
  \qquad\text{for all } x \in [0,1],
\]
where $B_n$ is the $n$-th Bernstein polynomial for $f$.

\smallskip
\noindent
\hyperref[prf:bernstein-approximation]{\textit{Go to proof.}}
\end{theorem}

\begin{remark*}[Standard quantified statement]
\[
  f \text{ cts on } [0,1],\;
  \varepsilon > 0
  \;\Rightarrow\;
  \exists n_\varepsilon \in \mathbb{N},\;
  \forall n \geq n_\varepsilon:\;
  \sup_{x \in [0,1]} |f(x) - B_n(x)| < \varepsilon.
\]
\end{remark*}

\begin{remark*}[Interpretation]
The Bernstein theorem is a constructive proof of the Weierstrass
Approximation Theorem. As $n \to \infty$, the Bernstein polynomials
converge uniformly to $f$ on $[0,1]$. This is the canonical
constructive proof: it builds the approximating polynomials
explicitly and the convergence proof is elementary, using only the
uniform continuity of $f$ and a variance estimate for the binomial
distribution. The Weierstrass theorem for general $[a,b]$ follows
by a linear change of variables.
\end{remark*}
